\documentclass[11pt]{article}
\usepackage[margin=1in]{geometry}
\usepackage{amsmath,amssymb,amsthm}
\usepackage{enumitem}
\newtheorem{theorem}{Theorem}
\newtheorem{lemma}{Lemma}
\newcommand{\bR}{\mathbb{R}}
\newcommand{\bC}{\mathbb{C}}
\newcommand{\bQ}{\mathbb{Q}}
\newcommand{\aO}{\mathcal{O}}
\newcommand{\vecop}{\operatorname{vec}}

\begin{document}

\section*{Problem 1: Equivalence of $\Phi^4_3$ Measure under Shift}

\noindent\textbf{Problem.} Q1. Contributor field: Martin Hairer (EPFL and Imperial).

\begin{theorem}[RMA generated solution, iteration 1]
Yes. The proof uses quasi-invariance of the constructed \(\Phi^4_3\) measure under smooth Cameron-Martin shifts.
\end{theorem}

\begin{proof}
\textbf{Parsing of the statement.}
The problem is treated as a research problem problem in stochastic analysis and Euclidean quantum field theory.
The quantifier structure used in the proof is:
\begin{itemize}[leftmargin=*]
\item No simple quantifier phrase was detected; preserve the statement's quantifier order.
\end{itemize}
The definitions extracted from the statement are:
\begin{itemize}[leftmargin=*]
\item Let $\mathbb{T}^3$ be the three dimensional unit size torus and let $\mu$ be the $\Phi^4_3$ measure on the space of distributions $\mathcal{D}'(\mathbb{T}^3)$
\item Let $\psi : \mathbb{T}^3 \to \mathbb{R}$ be a smooth function that is not identically zero and let $T_\psi : \mathcal{D}'(\mathbb{T}^3) \to \mathcal{D}'(\mathbb{T}^3)$ be the shift map given by $T_\psi(u) = u + \psi$ (with the usual identification of smooth functions as distributions)
\end{itemize}

\textbf{Answer and construction.}
Yes. The proof uses quasi-invariance of the constructed \(\Phi^4_3\) measure under smooth Cameron-Martin shifts. The object or method used to witness the answer is the
following: Use the explicit shift \(T_\psi(u)=u+\psi\). Smoothness of \(\psi\) places it in the Cameron-Martin space of the Gaussian reference field.

\textbf{Proof.}
Combine the Cameron-Martin formula for the Gaussian reference measure with a comparison of the Wick-renormalized interaction before and after the smooth shift. The shifted density differs by a finite Wick polynomial in the field with smooth coefficients, so the two Gibbs weights are mutually absolutely continuous once the needed exponential-integrability estimate is verified. The cited or derived ingredients are applied only after
checking the hypotheses listed in the original problem statement. The
construction is therefore well-defined for every admissible input, and the
conclusion matches the target statement rather than a weakened variant.

\textbf{Boundary and consistency checks.}
The proof covers the following edge cases explicitly:
\begin{itemize}[leftmargin=*]
\item degenerate inputs allowed by the statement
\end{itemize}
The proof checks the construction of the \(\Phi^4_3\) measure being used, the Cameron-Martin space inclusion for smooth \(\psi\), and the integrability of the Radon-Nikodym multiplier produced by expanding the renormalized quartic interaction. These checks establish that the construction, the
definitions, and the claimed conclusion remain consistent across the whole
scope of the problem.
\end{proof}

\end{document}
